\documentclass[10pt,a4paper]{amsart}
\usepackage[margin=19mm]{geometry}
\usepackage{amsmath,amssymb,mathtools}
\usepackage[scr=rsfs,cal=euler]{mathalfa}
\usepackage{xfrac}
\usepackage[unicode,hidelinks,bookmarks=false]{hyperref}
\newcommand{\ie}{i.e.\ }
\newcommand{\cf}{cf.\ }
\newcommand{\resp}{respectively\ }
\newcommand{\Subsection}{\subsection}
\providecommand{\nobreakdash}{\nobreak\hbox{-}}
\setlength{\emergencystretch}{2em}
\multlinegap0pt
\usepackage{mathtools}
\usepackage{amssymb}
\usepackage{braket}
\usepackage{xcolor}
\newtheorem{lemma}{Lemma}
\newcommand{\PSC}{\mathrm{PSC}}
\renewcommand{\geq}{\geqslant}
\renewcommand{\leq}{\leqslant}
\title{Binary code rate bounds via classical--quantum channels}
\author{Omar Alrabiah, Venkatesan Guruswami}
\date{Source: arXiv:2608.09347v1 (CC BY 4.0)}
\begin{document}
\maketitle

\noindent\emph{Source and licence.} Binary code rate bounds via classical--quantum channels. Version arXiv:2608.09347v1 (CC BY 4.0), distributed by arXiv under the Creative Commons Attribution 4.0 International licence, http://creativecommons.org/licenses/by/4.0/, as recorded in the arXiv licence metadata for this exact version and shown on the abstract page https://arxiv.org/abs/2608.09347v1. Full licence notice: \texttt{licenses/arxiv-2608.09347-CC-BY-4.0.txt}.\par\smallskip

\noindent\textit{Editorial scope.} Counted unit: the article's lemma lem:psc-sy-interval (source0.tex lines 2711--2719). The article's complete original proof of it is delivered with it (source0.tex lines 2721--2785, one \texttt{proof} environment of 65 lines). No mathematics is written, repaired or re-derived: the statement and the proof are the article's own lines, in the article's own notation, and no retained line is substituted or re-flowed.  Retained uncounted as the source-local context the proof needs: lines 2585--2595; lines 2699--2707.  Nothing was omitted from any retained span, so the package carries no omission record.  No retained line contains a citation, so the wrapper carries no bibliography and no bibliography entry.  No retained line contains a figure environment, an image-inclusion command or drawing code, so no image is delivered and the package declares no asset dependency.  Editorial formatting only, in addition: an AMS \texttt{amsart} preamble that restates the article's own declarations and macros used by the retained lines, namely the environments \texttt{lemma} on separate counters, together with the standard packages those lines need.  The retained environments print in this document as Lemma 1 in order of appearance, whereas the article prints the same items with the numbering of its own full document.  The complete native source of the article as distributed in the arXiv e-print for this exact version is bundled unchanged as \texttt{sources/207234/source0.tex}.
\begin{equation}
\label{eq:sy-ldpc-base}
B_3^{\mathrm{SY}}(\delta)
\coloneqq
\begin{cases}
\frac{2}{3}, & r_\delta \geq \frac{1}{4}, \\
-r_\delta\log{r_\delta}
-\frac{1-3r_\delta}{3}\log(1-3r_\delta),
& 0 < r_\delta < \frac{1}{4},
\end{cases}
\end{equation}

\begin{equation}
\label{eq:psc-ldpc-rate}
U_3^{\PSC}(\delta)
\coloneqq
\inf_{\substack{0 < p < \frac{1}{2}\\ e_3^{\mathrm{LDPC}}(p) < \delta}}
h\left(\frac{1}{2}-\sqrt{p(1-p)}\right)
=
h\left(\frac{1-\sqrt{\delta+\sqrt{\delta(2-3\delta)}}}{2}\right).
\end{equation}

\begin{lemma}[Shangguan--Yang interval separation]
\label{lem:psc-sy-interval}
The pointwise inequality
\begin{equation*}
U_3^{\PSC}(u)
< B_3^{\mathrm{SY}}(u)
\end{equation*}
holds for $\frac{1}{6} \leq u < \frac{1}{2}$.
\end{lemma}

\begin{proof}
For $u \in [\frac{1}{6},\frac{1}{2})$, set
\begin{equation*}
q
\coloneqq
\frac{1-\sqrt{u+\sqrt{u(2-3u)}}}{2},
\qquad
r
\coloneqq
\frac{1}{2}-\sqrt{u(1-u)}.
\end{equation*}
By \eqref{eq:psc-ldpc-rate}, $U_3^{\PSC}(u) = h(q)$, while $r$ is the parameter $r_\delta$ in \eqref{eq:sy-ldpc-base}. We compare both curves through $h(r)$.

The key estimate is
\begin{equation}
\label{eq:psc-w3-gain}
r-q > r^2.
\end{equation}
To prove it, put $z \coloneqq 1-2u \in (0,\frac{2}{3}]$. Then $r(1-r) = \frac{z^2}{4}$, so \eqref{eq:psc-w3-gain} is equivalent to $q < \frac{z^2}{4}$. After substituting $u = \frac{1-z}{2}$ and squaring twice, this inequality reduces to
\begin{equation*}
1+2z-3z^2
>
\left(1+z-2z^2+\frac{z^4}{2}\right)^2.
\end{equation*}
Both squarings preserve the inequality because all quantities involved are positive. Finally,
\begin{align*}
&1+2z-3z^2
-\left(1+z-2z^2+\frac{z^4}{2}\right)^2\\
&\qquad
=z^3\left(4-5z-z^2+2z^3-\frac{z^5}{4}\right).
\end{align*}
For $0 < z \leq \frac{2}{3}$, the factor in parentheses satisfies
\begin{equation*}
4-5z-z^2+2z^3-\frac{z^5}{4}
\geq
4-\frac{10}{3}-\frac{4}{9}-\frac{8}{243}
= \frac{46}{243} > 0.
\end{equation*}
This proves \eqref{eq:psc-w3-gain}.

The assumption $u \geq \frac{1}{6}$ gives
\begin{equation*}
r
\leq
\frac{3-\sqrt{5}}{6}
< \frac{1}{7}.
\end{equation*}
In particular, $B_3^{\mathrm{SY}}(u)$ lies on the entropy branch of \eqref{eq:sy-ldpc-base}. Since $h'(t) = \log\frac{1-t}{t}$ decreases on $(0,\frac{1}{2})$, \eqref{eq:psc-w3-gain} yields
\begin{equation}
\label{eq:hrq}
h(r)-h(q)
=\int_q^r h'(t)\,dt
\geq(r-q)\log\frac{1-r}{r}
>r^2\log{6}.
\end{equation}
On the other hand,
\begin{equation}
\label{eq:hSY}
h(r)-B_3^{\mathrm{SY}}(u)
=\int_0^r\log\frac{1-t}{1-3t}\,dt
\leq(\log{\mathrm e})\int_0^r\frac{2t}{1-3t}\,dt
<\frac{7}{4}(\log{\mathrm e})\,r^2.
\end{equation}
Here we used $\log(1+x) \leq (\log{\mathrm e})x$ and $1-3t > \frac{4}{7}$ on the integration interval. Since $\log{6} > \frac{7}{4}\log{\mathrm e}$, subtracting \eqref{eq:hSY} from \eqref{eq:hrq} gives $U_3^{\PSC}(u) < B_3^{\mathrm{SY}}(u)$.
\end{proof}

\end{document}
